% =======================================================================
% Relativistic Extension of Chapter V, Sections 49--50
% Combined Rotation and Magnetic Field Effects; Wave Propagation
%
% Conventions: (-,+,+,+) signature, c explicit, u^mu u_mu = -c^2
% Requires: rel_preamble.tex macros
% =======================================================================

\chapter{Relativistic Treatment of Combined Rotation and Magnetic Field
  Effects}\label{ch:rel-5-49-50}

This chapter extends Chandrasekhar's analysis of the combined effects of
rotation and magnetic field (Chapter~V, \S\S49--50) to the regime of
special relativity.  When either the Alfv\'en speed~$\vA$ or the
rotational velocity~$\Omega R$ becomes an appreciable fraction of the
speed of light~$c$, the Newtonian equations of Chapter~V must be replaced
by their fully relativistic counterparts.  The speed of light~$c$ is kept
explicit throughout, so that every classical limit can be recovered
transparently by taking $c\to\infty$.

% ======================================================================
\section{Like and contrary effects of rotation and magnetic field ---
  relativistic analysis}\label{sec:rel-49}
% ======================================================================

\subsection{Qualitative overview}

The qualitative interplay between rotation and magnetic field that
Chandrasekhar identifies in \S49 persists in the relativistic regime:
rotation inhibits convective instability by enforcing the
Taylor--Proudman constraint, a magnetic field does so through its
analogue (the tendency to freeze fluid elements along field lines), and
the two mechanisms nevertheless \emph{oppose} each other because the
magnetic field endows the fluid with an effective viscosity that
facilitates the onset of instability in the presence of rotation.

Relativistic corrections modify this picture quantitatively in three
principal ways.

\subsection{Relativistic inertia: the enthalpy--inertia coupling}

In special-relativistic hydrodynamics the role of the mass density
$\rho$ is assumed by the relativistic enthalpy density
\begin{equation}
  \enthalpy = \edensity + p,
  \label{eq:rel49-enthalpy}
\end{equation}
where $\edensity$ is the total energy density (including rest mass) and
$p$ the pressure.  The effective inertia per unit volume in the equation
of motion is $\enthalpy/c^{2}$, not~$\rho$.  Both rotation and the
magnetic field interact with this enhanced inertia:

\begin{enumerate}
\item \textbf{Rotation.}  The Coriolis acceleration in the co-rotating
  frame takes the covariant form
  \begin{equation}
    a^{\mu}_{\text{Cor}}
      = -\frac{2}{c^{2}}\,\epsilon^{\mu\nu\alpha\beta}\,
        u_{\nu}\,\Omega_{\alpha}\,u_{\beta}^{(\text{rot})},
    \label{eq:rel49-coriolis}
  \end{equation}
  where $\Omega_{\alpha}$ is the vorticity of the background rotation
  4-velocity.  The pre-factor $1/c^{2}$ (compared with the Newtonian
  Coriolis force $2\boldsymbol{\Omega}\times\mathbf{u}$) reflects the
  replacement $\rho\to \enthalpy/c^{2}$.

\item \textbf{Magnetic field.}  The magnetic tension that drives
  Alfv\'en waves now acts on $\enthalpy/c^{2}$ rather than $\rho$.
  The relativistic Alfv\'en speed is
  \begin{equation}
    \vA^{2}
      = \frac{b^{2}}{\enthalpy/c^{2} + b^{2}/c^{2}}\,c^{2}
      = \frac{b^{2}\,c^{2}}{\enthalpy + b^{2}},
    \label{eq:rel49-vA}
  \end{equation}
  where $b^{2}=b^{\mu}b_{\mu}$ is the invariant magnetic pressure.
  In the classical limit ($\enthalpy\to\rho c^{2}$,
  $b^{2}\ll\rho c^{2}$) this reduces to
  $\vA^{2}\to b^{2}/\rho = \mu H^{2}/(4\pi\rho)$ as required.

\item \textbf{Coupling term.}  The pressure contribution $p/c^{2}$
  to the inertia interacts simultaneously with \emph{both}~$\Omega$
  and~$\mathbf{B}$.  Denoting the fractional correction
  $\delta\equiv p/(\rho c^{2})$ (of order $v_{\text{th}}^{2}/c^{2}$
  for a non-degenerate gas), the combined rotational--magnetic
  frequency shifts acquire multiplicative factors $(1+\delta)^{-1}$
  that are absent in the Newtonian theory.  This is the mechanism
  through which the equation of state enters the rotational--magnetic
  interplay.
\end{enumerate}

\subsection{Relativistic Taylor and Chandrasekhar numbers}

The classical Taylor number $\mathrm{Ta}=4\Omega^{2}d^{4}/\nu^{2}$ and
Chandrasekhar number $Q=\mu H^{2}d^{2}/(\rho\nu\eta)$ generalise to
\begin{equation}
  \Tarel = \frac{4\Omega^{2}d^{4}}{\nu^{2}}
    \left(1 + \frac{p + b^{2}/2}{\rho c^{2}}\right)^{\!-2},
  \label{eq:rel49-Ta}
\end{equation}
\begin{equation}
  \Qrel = \frac{b^{2}\,d^{2}}{\enthalpy\,\nu\eta/c^{2}}
    = \frac{b^{2}\,c^{2}\,d^{2}}{\enthalpy\,\nu\eta},
  \label{eq:rel49-Q}
\end{equation}
where $\nu$ and $\eta$ are now kinematic viscosity and resistivity
measured in the fluid rest frame.  The ratio
\begin{equation}
  \frac{\Qrel}{\Tarel}
    = \frac{b^{2}\,c^{2}\,\nu}{4\Omega^{2}d^{2}\,\enthalpy\,\eta}
    \left(1 + \frac{p + b^{2}/2}{\rho c^{2}}\right)^{\!2}
  \label{eq:rel49-ratio}
\end{equation}
controls which effect dominates.  In the non-relativistic limit,
$\enthalpy\to\rho c^{2}$ and the parenthetical factor tends to unity,
recovering the classical result.

\subsection{Physical discussion of $p/c^{2}$ corrections}

The corrections $\order{p/(\rho c^{2})}$ modify the classical balance
of rotation and magnetic field in two concrete ways:
\begin{itemize}
\item \emph{Reduced effective rotation:} Since the inertia increases
  by a factor $(1+p/(\rho c^{2}))$, the effective angular velocity
  ``felt'' by perturbations is reduced:
  $\Omega_{\text{eff}}\approx\Omega/(1+p/(\rho c^{2}))$.  Rotation
  thus becomes \emph{less effective} at inhibiting convection in hot,
  high-pressure fluids.
\item \emph{Reduced Alfv\'en restoring force:} The magnetic tension per
  unit inertia likewise decreases, lowering $\vA$.  The Chandrasekhar
  number's stabilising influence is weakened for the same physical
  reason.
\item Because both effects are weakened by the \emph{same} factor, the
  ratio $\Qrel/\Tarel$ receives corrections only at
  $\order{p^{2}/(\rho c^{2})^{2}}$, so the qualitative competition
  between rotation and field is surprisingly robust against relativistic
  corrections at leading order.
\end{itemize}

\begin{relcorrection}
At leading relativistic order the classical result of \S49 is modified by
the replacement $\rho\to(\edensity+p)/c^{2}=\enthalpy/c^{2}$ in all
inertial terms.  Both rotation and magnetic field are weakened in absolute
terms, but their ratio is corrected only at second order in $v^{2}/c^{2}$.
\end{relcorrection}


% ======================================================================
\section{Hydromagnetic waves in a rotating fluid --- relativistic
  dispersion relation}\label{sec:rel-50}
% ======================================================================

We now extend the wave analysis of \S50 to special relativity.  The
background is a uniformly rotating perfect fluid permeated by a uniform
magnetic field, and we study small perturbations about this state.

\subsection{Setup and linearised equations}

We work in the co-rotating frame and adopt a locally Minkowski metric
$\eta^{\mu\nu}=\mathrm{diag}(-1,+1,+1,+1)$, valid when $\Omega R\ll c$
(i.e.\ the rotation is slow enough that the metric is only weakly
non-inertial).  The background magnetic 4-vector $b_{0}^{\mu}$ satisfies
$b_{0}^{\mu}u^{(0)}_{\mu}=0$ and $b_{0}^{2}=b_{0}^{\mu}b_{0\mu}$.

Linearising the relativistic MHD equations about this background, using
the ideal MHD condition $F^{\mu\nu}u_{\nu}=0$ and
$u^{\mu}u_{\mu}=-c^{2}$, and seeking plane-wave solutions
$\propto\exp[i(\omega t - \mathbf{k}\cdot\mathbf{x})]$, one arrives at
\begin{equation}
  \omega\,\delta u_{i}
    - \frac{b_{0}^{2}}{\enthalpy/c^{2}+b_{0}^{2}/c^{2}}
      \,(k_{j}\hat{b}_{j})\,\delta b_{i}
    = -k_{i}\,\delta\varpi
      + 2i\,\epsilon_{ijk}\,\delta u_{j}\,\Omega_{k},
  \label{eq:rel50-mom}
\end{equation}
\begin{equation}
  \omega\,\delta b_{i}
    = (k_{j}\hat{b}_{j})\,|\mathbf{b}_{0}|\,\delta u_{i},
  \label{eq:rel50-ind}
\end{equation}
\begin{equation}
  k_{i}\,\delta u_{i} = 0,
  \qquad
  k_{i}\,\delta b_{i} = 0,
  \label{eq:rel50-div}
\end{equation}
where $\hat{b}_{j}=b_{0j}/|\mathbf{b}_{0}|$ is the unit vector along
the background field direction, and
$\delta\varpi$ is the perturbation of the total (fluid + magnetic)
pressure.  The key relativistic modification in~\eqref{eq:rel50-mom} is
the replacement
\begin{equation}
  \frac{\mu}{4\pi\rho}
  \;\longrightarrow\;
  \frac{b_{0}^{2}}{\enthalpy/c^{2}+b_{0}^{2}/c^{2}}
  = \frac{\vA^{2}}{c^{2}}\,\frac{\enthalpy+b_{0}^{2}}{b_{0}^{2}}\cdot
    \frac{b_{0}^{2}}{1} \cdot\frac{c^{2}}{\enthalpy+b_{0}^{2}}
  = \vA^{2},
  \label{eq:rel50-replacement}
\end{equation}
in terms of the relativistic Alfv\'en speed~\eqref{eq:rel49-vA}.

\subsection{Elimination of the magnetic perturbation}

Exactly as in the classical analysis, we eliminate $\delta b_{i}$ using
\eqref{eq:rel50-ind} to obtain
\begin{equation}
  n_{\mathrm{rel}}\,\delta u_{i}
    + k_{i}\,\delta\varpi
    + 2i\,\epsilon_{ijk}\,\delta u_{j}\,\Omega_{k} = 0,
  \label{eq:rel50-reduced}
\end{equation}
where the relativistic analogue of Chandrasekhar's $n$ is
\begin{equation}
  n_{\mathrm{rel}}
    = \frac{1}{\omega}
      \bigl(\omega^{2} - \omega_{A,\mathrm{rel}}^{2}\bigr),
  \label{eq:rel50-n}
\end{equation}
with the relativistic Alfv\'en frequency
\begin{equation}
  \omega_{A,\mathrm{rel}}
    = \vA\,(k_{j}\hat{b}_{j})
    = \vA\,k\cos\phi,
  \label{eq:rel50-omegaA}
\end{equation}
where $\phi$ is the angle between the wave vector and the magnetic field
direction.  Since equation~\eqref{eq:rel50-reduced} has exactly the same
algebraic structure as the classical equation~(10) of Chapter~V, the same
solution procedure applies.

\subsection{Full relativistic dispersion relation}

Choosing coordinates so that the $z$-axis lies along $\mathbf{k}$ and the
$(x,z)$-plane contains $\boldsymbol{\Omega}$, the transversality
condition~\eqref{eq:rel50-div} gives $\delta u_{z}=0$.  The resulting
$2\times 2$ system for $(\delta u_{x},\delta u_{y})$ yields
\begin{equation}
  n_{\mathrm{rel}} = \pm\,2\Omega\cos\vartheta,
  \label{eq:rel50-n-result}
\end{equation}
where $\vartheta$ is the angle between $\mathbf{k}$ and
$\boldsymbol{\Omega}$.  Substituting the
definition~\eqref{eq:rel50-n} gives the \textbf{relativistic dispersion
relation}:
\begin{equation}
  \boxed{%
    \omega^{2}
    \mp (2\Omega\cos\vartheta)\,\omega
    - \vA^{2}\,k^{2}\cos^{2}\!\phi
    = 0.
  }
  \label{eq:rel50-dispersion}
\end{equation}
This is formally identical to the classical equation~(16) of \S50, with
the single but crucial replacement
$\mu H_{j}k_{j}/(4\pi\rho)^{1/2}\to\vA\,k\cos\phi$, where $\vA$ is the
relativistic Alfv\'en speed~\eqref{eq:rel49-vA}.

The solutions are
\begin{equation}
  \omega = \pm\Bigl[
    \Omega\cos\vartheta
    \pm\sqrt{\Omega^{2}\cos^{2}\!\vartheta
      + \vA^{2}\,k^{2}\cos^{2}\!\phi}
  \,\Bigr].
  \label{eq:rel50-roots}
\end{equation}

\subsubsection{Including compressibility and thermal effects}

When the fluid is compressible and the sound speed $\cs$ is finite, the
longitudinal mode couples to the transverse modes.  The full
$3\times 3$ dispersion relation becomes
\begin{equation}
  \boxed{%
    \bigl(\omega^{2} - \cs^{2}k^{2}\bigr)
    \bigl(\omega^{2} - \vA^{2}k^{2}\cos^{2}\!\phi\bigr)
    - 4\Omega^{2}\cos^{2}\!\vartheta\,\omega^{2}
    + \cs^{2}k^{2}\,\vA^{2}k^{2}\cos^{2}\!\phi\,\sin^{2}\!\alpha
    = 0,
  }
  \label{eq:rel50-full-dispersion}
\end{equation}
where $\alpha$ is the angle between $\mathbf{k}$ and the magnetic field,
and both $\cs$ and $\vA$ are the \emph{relativistic} sound and Alfv\'en
speeds satisfying $\cs^{2}<c^{2}$ and $\vA^{2}<c^{2}$.

This is a cubic in $\omega^{2}$; it reduces to the classical
Chandrasekhar result when $\cs,\vA\ll c$ and $\Omega R\ll c$.

For general propagation, $\omega=\omega(k,\vartheta,\phi,\Omega,\vA,
\cs,c)$ with $c$ entering through the relativistic definitions of
$\cs$ and~$\vA$:
\begin{equation}
  \cs^{2} = \left(\frac{\partial p}{\partial\edensity}\right)_{\!s}
  \leq c^{2},
  \qquad
  \vA^{2} = \frac{b^{2}\,c^{2}}{\enthalpy + b^{2}} < c^{2}.
  \label{eq:rel50-speeds}
\end{equation}

\subsection{Limiting cases and mode identification}

\paragraph{(a) No rotation ($\Omega=0$).}
Equation~\eqref{eq:rel50-full-dispersion} factorises into the standard
relativistic MHD modes:
\begin{itemize}
\item \emph{Alfv\'en mode:}
  $\omega^{2}=\vA^{2}\,k^{2}\cos^{2}\!\phi$.
\item \emph{Fast/slow magnetosonic modes:}
  $\omega^{2}=\tfrac{1}{2}k^{2}
    \bigl[(\cs^{2}+\vA^{2})
      \pm\sqrt{(\cs^{2}+\vA^{2})^{2}
        -4\cs^{2}\vA^{2}\cos^{2}\!\phi}\,\bigr]$.
\end{itemize}
The fast magnetosonic speed is bounded by
$\vf^{2}=\cs^{2}+\vA^{2}-\cs^{2}\vA^{2}/c^{2}<c^{2}$, confirming
sub-luminal propagation (see \S\ref{subsec:causality} below).

\paragraph{(b) No magnetic field ($\vA=0$).}
The dispersion relation becomes
\begin{equation}
  \omega^{2}\bigl(\omega^{2}-\cs^{2}k^{2}\bigr)
    = 4\Omega^{2}\cos^{2}\!\vartheta\,\omega^{2},
\end{equation}
yielding
$\omega^{2}=\cs^{2}k^{2}+4\Omega^{2}\cos^{2}\!\vartheta$ (the
relativistic inertial--acoustic mode) and $\omega=0$ (the geostrophic
mode), recovering the relativistic rotating-fluid result.

\paragraph{(c) Incompressible limit ($\cs\to\infty$, capped at $c$).}
In the relativistic framework one cannot literally send $\cs\to\infty$;
instead one takes $\cs\to c$.  The longitudinal mode decouples at the
causal limit and the remaining transverse dispersion
relation~\eqref{eq:rel50-dispersion} is recovered.

\paragraph{(d) Non-relativistic limit ($c\to\infty$).}
In this limit $\vA^{2}\to\mu H^{2}/(4\pi\rho)$, $\cs^{2}\to(\gamma
p/\rho)$, and all equations reduce to those of \S50 in Chapter~V.

\paragraph{(e) Aligned case ($\mathbf{B}\parallel\boldsymbol{\Omega}$,
$\phi=\vartheta$).}
The magnetic and rotational axes coincide.  Writing
$\theta\equiv\vartheta=\phi$, the incompressible dispersion relation
becomes
\begin{equation}
  \omega = \pm\bigl[\Omega\cos\theta
    \pm\sqrt{\Omega^{2}\cos^{2}\!\theta + \vA^{2}k^{2}\cos^{2}\!\theta}
    \,\bigr],
\end{equation}
which can be written as
\begin{equation}
  \omega = \cos\theta\,\bigl[\pm\Omega
    \pm\sqrt{\Omega^{2}+\vA^{2}k^{2}}\,\bigr].
  \label{eq:rel50-aligned}
\end{equation}
The four branches are: two magneto-inertial waves (rotation dominates for
small $k$) and two magneto-Alfv\'en waves (field dominates for large $k$).

\subsection{Phase and group velocities}

The phase velocity for the incompressible modes is
\begin{equation}
  V_{\text{ph}} = \frac{\omega}{k}
    = \pm\frac{1}{k}\Bigl[\Omega\cos\vartheta
      \pm\sqrt{\Omega^{2}\cos^{2}\!\vartheta
        + \vA^{2}k^{2}\cos^{2}\!\phi}\,\Bigr],
  \label{eq:rel50-phase}
\end{equation}
and the group velocity is
\begin{equation}
  \frac{\partial\omega}{\partial\mathbf{k}}
    = \pm\left\{
      \frac{\mathbf{k}\times(\boldsymbol{\Omega}\times\mathbf{k})}{k^{3}}
      \pm\frac{\Omega\cos\vartheta\,
        [\mathbf{k}\times(\boldsymbol{\Omega}\times\mathbf{k})]\,k^{-3}
        + \vA^{2}\,(k_{j}\hat{b}_{j})\,\hat{\mathbf{b}}}%
      {\sqrt{\Omega^{2}\cos^{2}\!\vartheta
        + \vA^{2}k^{2}\cos^{2}\!\phi}}
    \right\}.
  \label{eq:rel50-group}
\end{equation}
This reduces to the classical expression~(20) of \S50 when
$\vA$ is replaced by its Newtonian value.

\subsection{Causality verification}\label{subsec:causality}

\begin{causalitycheck}
We verify that the group velocity $|\partial\omega/\partial\mathbf{k}|$
does not exceed $c$ for any wave branch.

\paragraph{Incompressible transverse modes.}
From~\eqref{eq:rel50-roots}, the frequency satisfies
$|\omega|\leq|\Omega\cos\vartheta|+
\sqrt{\Omega^{2}\cos^{2}\!\vartheta+\vA^{2}k^{2}\cos^{2}\!\phi}$.
For large $k$ the dominant term is $\vA k|\cos\phi|$, giving a group
speed asymptotic to $\vA|\cos\phi|\leq\vA<c$ (since $\vA<c$ by
construction, equation~\eqref{eq:rel49-vA}).  The rotation
contribution~$\Omega\cos\vartheta$ is $k$-independent and so contributes
a term $\sim 1/k$ to the group velocity at large~$k$, which vanishes.
Hence:
\begin{equation}
  \bigl|\mathbf{v}_{g}\bigr|
    \leq \vA < c
  \qquad\text{for the incompressible modes.}
  \label{eq:rel50-causality-incomp}
\end{equation}

\paragraph{Compressible modes (full dispersion).}
The fast magnetosonic speed in the rotating case is bounded by the
non-rotating fast speed:
\begin{equation}
  v_{g,\max}^{2}
    \leq \vf^{2}
    = \cs^{2} + \vA^{2} - \frac{\cs^{2}\vA^{2}}{c^{2}}
    < c^{2},
  \label{eq:rel50-causality-comp}
\end{equation}
since $\cs^{2}<c^{2}$ and $\vA^{2}<c^{2}$.  The inequality
$\vf^{2}<c^{2}$ follows from
\begin{equation}
  c^{2}-\vf^{2}
    = c^{2}-\cs^{2}-\vA^{2}+\frac{\cs^{2}\vA^{2}}{c^{2}}
    = \frac{(c^{2}-\cs^{2})(c^{2}-\vA^{2})}{c^{2}} > 0.
  \label{eq:rel50-causality-proof}
\end{equation}
Rotation modifies the mode frequencies but does not increase the
maximum group speed above $\vf$, because the Coriolis force does no work
(it is perpendicular to the velocity) and therefore cannot increase the
energy propagation speed.

\paragraph{Slow mode.}
The slow magnetosonic group speed satisfies
$v_{g,\mathrm{slow}}^{2}\leq\vs^{2}=\cs^{2}\vA^{2}/\vf^{2}<\vf^{2}<c^{2}$.

\medskip
\noindent\textbf{Conclusion:} All wave branches propagate sub-luminally.
The combined rotation + magnetic field system is causal in the
relativistic regime.
\end{causalitycheck}

\subsection{Summary of relativistic modifications}

\begin{relcorrection}
The relativistic extension of \S\S49--50 introduces three principal
modifications to the classical theory:
\begin{enumerate}
\item The mass density $\rho$ is replaced everywhere by the relativistic
  inertia $\enthalpy/c^{2}=(\edensity+p)/c^{2}$.  This weakens both
  rotational and magnetic stabilisation in hot, high-pressure fluids.
\item The Alfv\'en speed becomes the relativistic expression
  $\vA^{2}=b^{2}c^{2}/(\enthalpy+b^{2})$, which is automatically bounded
  by~$c$.
\item The sound speed satisfies $\cs^{2}\leq c^{2}$; the fast
  magnetosonic speed satisfies
  $\vf^{2}=\cs^{2}+\vA^{2}-\cs^{2}\vA^{2}/c^{2}<c^{2}$; and all group
  velocities are bounded by~$c$.  Causality is guaranteed for all
  wave branches.
\end{enumerate}
In the limit $c\to\infty$ every result reduces to the corresponding
expression in Chandrasekhar's Chapter~V, \S\S49--50.
\end{relcorrection}
